% 教材：Shankar, Principles of Quantum Mechanics
% 精确范围：第 1 章第 1.1 节，印刷页 pp. 1--7，PDF pp. 15--21
% 人工核验：已完成讲解、问答和教材练习 1.1.1--1.1.5

\section{第 1.1 节：线性向量空间基础}
\label{section:QM-C01-S01}

\studymeta
  {QM-C01-S01}
  {2026-08-17}
  {Shankar, \textit{Principles of Quantum Mechanics}}
  {第 1 章 \S 1.1；印刷页 pp.~1--7；PDF pp.~15--21}
  {已确认（讲解、问答与教材练习核验）}

\subsection{本节主线}

线性向量空间把箭头、矩阵、函数等外形不同的对象统一起来。统一它们的不是长度和方向，而是两种运算：向量加法与标量数乘。公理确定哪些推理对所有向量空间都成立；线性无关、维数和基进一步说明如何用一组坐标唯一表示抽象向量。

\notetopic{QM-C01-S01-T01}{向量空间、标量域与公理}

\keyidea{“向量”首先是向量空间中的元素，不一定是箭头或一列数。判断一个集合是否为向量空间，要先说明允许的标量以及向量加法、标量数乘的规则。}

\begin{definition}
设 $V$ 是一个集合，标量来自数域 $\mathbb{F}$。若 $V$ 上定义了向量加法与标量数乘，并满足下列条件，则称 $V$ 为 $\mathbb{F}$ 上的线性向量空间：
\begin{enumerate}
  \item 对任意 $\ket{V},\ket{W}\in V$，有 $\ket{V}+\ket{W}\in V$；对任意 $a\in\mathbb{F}$，有 $a\ket{V}\in V$；
  \item $\ket{V}+\ket{W}=\ket{W}+\ket{V}$；
  \item $\ket{V}+(\ket{W}+\ket{Z})=(\ket{V}+\ket{W})+\ket{Z}$；
  \item 存在唯一的零向量 $\ket{0}$，使 $\ket{V}+\ket{0}=\ket{V}$；
  \item 每个 $\ket{V}$ 都有加法逆元 $-\ket{V}$，使 $\ket{V}+(-\ket{V})=\ket{0}$；
  \item $a(\ket{V}+\ket{W})=a\ket{V}+a\ket{W}$；
  \item $(a+b)\ket{V}=a\ket{V}+b\ket{V}$；
  \item $a(b\ket{V})=(ab)\ket{V}$，且 $1\ket{V}=\ket{V}$。
\end{enumerate}
\end{definition}

数域决定允许使用的标量。$\mathbb{F}=\real$ 时称为实向量空间，$\mathbb{F}=\complex$ 时称为复向量空间；“实”或“复”修饰的是标量，不要求向量对象本身看起来像实数或复数。

\begin{remark}
教材此处列出的公理没有单独写出 $1\ket{V}=\ket{V}$。本笔记按标准线性代数定义将标量单位元作用写明；后续涉及坐标和基时默认采用这一标准条件。
\end{remark}

由公理可以推出：零向量唯一，$0\ket{V}=\ket{0}$，$(-1)\ket{V}=-\ket{V}$，并且每个向量的加法逆元唯一。完整证明见\relatedexercise{QM-C01-S01-E01}{教材练习 1.1.1}。

\notetopic{QM-C01-S01-T02}{典型例子、封闭性与齐次约束}

箭头在通常的首尾相接加法和实数伸缩下构成实向量空间，但只取 $z$ 分量为正的箭头不构成向量空间：零向量不在集合中，负数数乘和加法逆元也会离开集合。

所有 $2\times 2$ 实矩阵在逐项相加与实数数乘下构成向量空间。零向量是零矩阵，加法逆元由每个矩阵元取负得到。矩阵在这里是“向量空间的元素”；矩阵乘法不是定义该向量空间所必需的运算。

区间 $0\leq x\leq L$ 上的函数在逐点相加和数乘下也可构成向量空间。例如，满足
\[
  f(0)=f(L)=0
\]
或
\[
  f(0)=f(L)
\]
的函数集合都对加法和数乘封闭，并包含零函数与加法逆函数。相反，满足 $f(0)=4$ 的函数集合不包含零函数，且一般不对加法、数乘和取负封闭。

\dialoguesupplement{形如 $L(f)=0$ 的齐次线性约束通常定义一个向量子空间；形如 $L(f)=c\neq0$ 的非齐次约束通常不构成向量空间。这是快速判断规则，最终仍须检查零向量和两种封闭性。}

\notetopic{QM-C01-S01-T03}{线性组合、线性相关与线性无关}

\begin{definition}
对向量组 $\ket{1},\ldots,\ket{n}$，表达式
\[
  \sum_{i=1}^{n}a_i\ket{i}=\ket{0}
\]
称为一条线性关系。若它只有平凡解 $a_1=\cdots=a_n=0$，该向量组线性无关；若存在不全为零的系数使等式成立，该向量组线性相关。
\end{definition}

若一条非平凡线性关系中 $a_k\neq0$，则
\[
  \ket{k}=-\sum_{i\neq k}\frac{a_i}{a_k}\ket{i}.
\]
因此线性相关等价于至少有一个向量可以由其余向量线性表示。证明线性相关只需找到一个非平凡关系；证明线性无关则必须排除所有非平凡关系。

平面内两个不平行的非零箭头线性无关，而任意三个平面箭头都线性相关。教材练习中的矩阵关系
\[
  \ket{1}-2\ket{2}=\ket{3}
\]
也直接给出非平凡线性关系，见\relatedexercise{QM-C01-S01-E04}{教材练习 1.1.4}。

\notetopic{QM-C01-S01-T04}{维数、基与展开定理}

\begin{definition}
一个向量空间能够容纳的线性无关向量的最大数目称为该空间的维数。若最大数目为 $n$，则该空间是 $n$ 维的。
\end{definition}

维数依赖标量域。例如 $\complex$ 作为复向量空间的维数为 $1$，基可取 $\{1\}$；作为实向量空间的维数为 $2$，基可取 $\{1,\ii\}$。

\begin{theorem}
在 $n$ 维向量空间中，任意 $n$ 个线性无关向量都能张成整个空间。
\end{theorem}

\textbf{证明。}设 $\ket{1},\ldots,\ket{n}$ 线性无关。若存在 $\ket{V}$ 不能由它们线性表示，则 $\ket{1},\ldots,\ket{n},\ket{V}$ 仍然线性无关：在线性关系
\[
  a\ket{V}+\sum_{i=1}^{n}b_i\ket{i}=\ket{0}
\]
中，若 $a\neq0$，就能解出 $\ket{V}$ 是其余向量的线性组合，与假设矛盾；故 $a=0$，再由原向量组线性无关得到全部 $b_i=0$。这将产生 $n+1$ 个线性无关向量，与空间为 $n$ 维矛盾。因此任意 $\ket{V}$ 都在这 $n$ 个向量的张成空间中。\hfill$\square$

\begin{definition}
一组既线性无关又张成整个空间的向量称为一组基。在已知空间为 $n$ 维时，这等价于取 $n$ 个线性无关向量。
\end{definition}

于是任意向量都可以写为
\begin{equation}
  \ket{V}=\sum_{i=1}^{n}v_i\ket{i},
  \label{eq:QM-C01-S01-basis-expansion}
\end{equation}
其中 $\{\ket{i}\}_{i=1}^{n}$ 是一组基，$v_i$ 称为 $\ket{V}$ 在该基下的分量或坐标。

\notetopic{QM-C01-S01-T05}{坐标唯一性与换基}

\begin{theorem}
同一个向量在给定的一组基下具有唯一的坐标。
\end{theorem}

\textbf{证明。}若
\[
  \ket{V}=\sum_i v_i\ket{i}=\sum_i v_i'\ket{i},
\]
则相减得到
\[
  \sum_i(v_i-v_i')\ket{i}=\ket{0}.
\]
基向量线性无关，所以每个系数差都必须为零，即 $v_i=v_i'$。\hfill$\square$

唯一性中的“给定一组基”不可省略。更换基以后，坐标通常改变，但抽象向量本身不变。例如若
\[
  \ket{1'}=\ket{1}+\ket{2},\qquad
  \ket{2'}=\ket{1}-\ket{2},
\]
则同一个向量可以写成
\[
  \ket{V}=3\ket{1}-2\ket{2}
  =\frac{1}{2}\ket{1'}+\frac{5}{2}\ket{2'}.
\]
第二个表达式可通过展开新基并比较 $\ket{1},\ket{2}$ 的系数验证。

若
\[
  \ket{V}=\sum_i v_i\ket{i},\qquad
  \ket{W}=\sum_i w_i\ket{i},
\]
则
\[
  \ket{V}+\ket{W}=\sum_i(v_i+w_i)\ket{i},\qquad
  a\ket{V}=\sum_i(av_i)\ket{i}.
\]
所以选定基以后，抽象的向量运算转化为对应分量的普通运算。

\notetopic{QM-C01-S01-T06}{矩阵与函数的坐标表示：对话补充}

任意实 $2\times2$ 矩阵都可唯一写成
\[
  \begin{pmatrix}a&b\\c&d\end{pmatrix}
  =a\begin{pmatrix}1&0\\0&0\end{pmatrix}
  +b\begin{pmatrix}0&1\\0&0\end{pmatrix}
  +c\begin{pmatrix}0&0\\1&0\end{pmatrix}
  +d\begin{pmatrix}0&0\\0&1\end{pmatrix}.
\]
因此 $M_{2\times2}(\real)$ 与 $\real^4$ 作为向量空间同构，矩阵的坐标可写成 $(a,b,c,d)^{\mathsf T}$。保留矩阵形式的好处是能够直接看出行列结构，并自然表达矩阵乘法、线性映射复合和特征值等额外结构。

有限维函数空间同样可以坐标化。例如
\[
  P_2=\{a+bx+cx^2:a,b,c\in\real\}
\]
在基 $\{1,x,x^2\}$ 下与 $\real^3$ 同构。一般函数空间可能是无限维的；只记录有限个采样值通常会丢失信息，除非事先限制函数类别。保留函数形式则便于表达连续变量、微分、积分和边界条件。

\begin{remark}
本知识点是学习对话中的补充解释，不是教材第 1.1 节的独立定理。其作用是连接“抽象向量”与熟悉的坐标列向量；无限维空间的严格坐标展开需要后续的内积与完备性概念。
\end{remark}

\subsection{一分钟回顾}

向量空间由标量域、向量加法、标量数乘和相应公理确定。线性无关意味着不存在非平凡的零向量线性关系；维数是线性无关向量数目的上限；基同时具备线性无关和张成性质。给定基后，每个向量具有唯一坐标，向量加法和标量数乘可逐分量进行。矩阵和函数都可以是向量，列数组只是选定基之后的坐标表达。
